\chapter{The Eudaimonia Society}

\section{Initial Exploration}

\subsection{General Thoughts \& Scope}
Originally, I had begun this exercise by thinking of, if I were to create an institution like the Freemasonry, what would it look like. What doctrine would I lay at the foundation, traditions I would create, values to instill, how to even go about it, and much more. This is an attempt to write it out. One more pragmatic train of thought, I would likely use the general idea in creating a club at school, very similar to my previous SVA, but more evolved. This club would likely be more constrained, focusing on intellectual virtues and their relation to research rather than the broad idea of a society based upon my fundamental values and then using them to grow to be something more.

I will break this up into Moral development, intellectual development, social cohesion, and social impact.

\subsubsection{Moral Development}
The core of the entire project is to take the fundamental approach towards perfection to be the absolute goal. To take all things seriously, never do things half-heartedly, and don't mistakenly take irony to that which is sacred.

That we are all within our own personal design and can take control of our lives and personalities through the continuous cultivation of habits.

On a more direct statement; take accountability for all things(even that which many would consider beyond you), never do anything anti-social(in the psychological sense) even in the most slight way, value truth, be accurate and develop yourself epistemically and intellectually, pursue the development of your community in your own image and values, develop clearly defensible and defined values that are not vague bromides, hold yourself to your standards rigorously, think in the long term for your community, only acception to the anti-social rule would be when it conflicts with written values, and be fully conscious at all times, fully aware of your actions and their relation to your moral identity.

I personally find virtue ethics to be my way of thinking, so this may be a core tenet or may be too constricting.

\subsubsection{Intellectual Development}
To understand things at their most fundamental level. To understand epistemology, statistical thinking, critical thinking, interdisciplinary thinking, metaphysics, and have a fully simulated world-view rather than implicitly taking on assumptions. This is the core tenent, not to take on others assumptions passed through pop culture.

\subsubsection{Social Cohesion}
Generally thinking a somewhat brotherhood of intellectuals or in close enough fields. To discuss and understand each other and build one another up.

Beyond that many questions arise of what structure would best work for this.

\subsubsection{Social Impact}

Any general social impact will do

\subsection{First layer of Refinement}
The first questions I want to ask consist of:
\begin{itemize}
  \item Mission Statement
  \item Structure
  \item Symbols(if any)

\end{itemize}
